\documentclass[11pt,a4paper]{article}
\usepackage[margin=2.5cm]{geometry}
\usepackage{amsmath,amssymb}
\usepackage{graphicx}
\usepackage{tcolorbox}
\usepackage{listings}
\usepackage{xcolor}
\usepackage{hyperref}
\usepackage{tikz}
\usetikzlibrary{shapes,arrows,positioning,fit}

% Color definitions
\definecolor{mlblue}{RGB}{0,102,204}
\definecolor{mlpurple}{RGB}{51,51,178}
\definecolor{mlorange}{RGB}{255,127,14}
\definecolor{mlgreen}{RGB}{44,160,44}

% Key concept box
\newtcolorbox{keybox}[1]{
  colback=mlblue!5!white,
  colframe=mlblue,
  fonttitle=\bfseries,
  title=#1
}

% Example box
\newtcolorbox{examplebox}[1]{
  colback=mlgreen!5!white,
  colframe=mlgreen,
  fonttitle=\bfseries,
  title=#1
}

% Warning box
\newtcolorbox{warningbox}[1]{
  colback=mlorange!5!white,
  colframe=mlorange,
  fonttitle=\bfseries,
  title=#1
}

% Code listing style
\lstset{
  basicstyle=\ttfamily\small,
  backgroundcolor=\color{gray!10},
  frame=single,
  numbers=left,
  numberstyle=\tiny\color{gray},
  breaklines=true,
  captionpos=b
}

% Solidity style
\lstdefinelanguage{Solidity}{
  keywords={contract, function, public, private, view, pure, payable, returns, address, uint, bool, string, mapping, struct, event, modifier, require, assert, revert, if, else, for, while, return, emit, constructor, memory, storage, calldata},
  comment=[l]{//},
  morecomment=[s]{/*}{*/},
  morestring=[b]",
  sensitive=true
}

\title{\textbf{Lecture Notes: Lesson 3}\\Ethereum and Smart Contracts}
\author{Cryptocurrency Course}
\date{}

\begin{document}

\maketitle

\tableofcontents
\newpage

\section{Introduction to Ethereum}

\subsection{Historical Background}

Ethereum was proposed by Vitalik Buterin in late 2013 and launched in July 2015. It extends Bitcoin's blockchain concept to support arbitrary computation through smart contracts.

\begin{keybox}{Ethereum's Vision}
While Bitcoin is \textbf{digital gold} (store of value), Ethereum is a \textbf{world computer} enabling:
\begin{itemize}
  \item Decentralized applications (DApps)
  \item Programmable money and assets
  \item Trustless agreements (smart contracts)
  \item Decentralized finance (DeFi)
  \item Non-fungible tokens (NFTs)
\end{itemize}
\end{keybox}

\subsection{Key Differences from Bitcoin}

\begin{center}
\begin{tabular}{|l|l|l|}
\hline
\textbf{Feature} & \textbf{Bitcoin} & \textbf{Ethereum} \\
\hline
Purpose & Digital currency & World computer \\
Scripting & Limited (Bitcoin Script) & Turing-complete (EVM) \\
Block Time & ~10 minutes & ~12 seconds \\
Consensus (2024) & Proof-of-Work & Proof-of-Stake \\
Supply Cap & 21 million BTC & No fixed cap \\
Account Model & UTXO & Account-based \\
\hline
\end{tabular}
\end{center}

\section{Ethereum Virtual Machine (EVM)}

\subsection{Architecture Overview}

The EVM is a stack-based virtual machine that executes smart contract bytecode. It provides a sandboxed execution environment where code runs deterministically across all nodes.

\subsubsection{Key Characteristics}

\begin{itemize}
  \item \textbf{Turing-Complete}: Can express any computation (with gas limits)
  \item \textbf{Deterministic}: Same input always produces same output
  \item \textbf{Isolated}: Contracts run in sandbox, can't access file system/network
  \item \textbf{Persistent}: State stored on blockchain between executions
\end{itemize}

\subsection{EVM Components}

\subsubsection{Stack}

\begin{itemize}
  \item 1024 element maximum depth
  \item Each element is 256 bits (32 bytes)
  \item All operations work on stack items
  \item LIFO (Last In, First Out) structure
\end{itemize}

\textbf{Example Operations}:
\begin{verbatim}
PUSH1 0x05    // Push 5 onto stack
PUSH1 0x03    // Push 3 onto stack
ADD           // Pop two items, push sum (8)
PUSH1 0x02    // Push 2 onto stack
MUL           // Pop two items, push product (16)
\end{verbatim}

\subsubsection{Memory}

\begin{itemize}
  \item Volatile (cleared between transactions)
  \item Byte-addressable array
  \item Expands dynamically (costs gas)
  \item Used for temporary data during execution
\end{itemize}

\subsubsection{Storage}

\begin{itemize}
  \item Persistent key-value store
  \item Each contract has own storage
  \item 256-bit keys and values
  \item Expensive to use (high gas cost)
  \item Survives between transactions
\end{itemize}

\begin{keybox}{Storage Cost}
Storage is the most expensive operation:
\begin{itemize}
  \item \textbf{SSTORE} (write): 20,000 gas (new) or 5,000 gas (update)
  \item \textbf{SLOAD} (read): 200 gas
  \item \textbf{Memory}: ~3 gas per word
\end{itemize}

Optimization strategy: Minimize storage writes, use events for logging.
\end{keybox}

\subsection{Account Model}

Ethereum uses an account-based model (unlike Bitcoin's UTXO model).

\subsubsection{Externally Owned Accounts (EOA)}

Controlled by private keys:
\begin{itemize}
  \item Address derived from public key (last 20 bytes of Keccak-256 hash)
  \item Balance in wei ($1$ ETH $= 10^{18}$ wei)
  \item Nonce (transaction counter)
  \item No code
\end{itemize}

\subsubsection{Contract Accounts}

Controlled by code:
\begin{itemize}
  \item Address derived from creator address + nonce
  \item Balance in wei
  \item Nonce (creation counter)
  \item Code hash (immutable bytecode)
  \item Storage root (Merkle tree of state)
\end{itemize}

\subsection{State Transition Function}

Ethereum's state machine transitions are deterministic:
\[
\sigma_{t+1} = \Upsilon(\sigma_t, T)
\]

where:
\begin{itemize}
  \item $\sigma_t$ is the world state at block $t$
  \item $T$ is a transaction
  \item $\Upsilon$ is the state transition function
  \item $\sigma_{t+1}$ is the resulting new state
\end{itemize}

\section{Gas Mechanism}

\subsection{Purpose of Gas}

Gas serves three critical functions:

\begin{enumerate}
  \item \textbf{Metering}: Limits computation to prevent infinite loops
  \item \textbf{Pricing}: Aligns costs with resource usage (computation, storage)
  \item \textbf{Fee Market}: Prioritizes transactions during congestion
\end{enumerate}

\begin{keybox}{Gas Economics}
\begin{align*}
\text{Transaction Cost} &= \text{Gas Used} \times \text{Gas Price} \\
\text{Max Fee} &= \text{Gas Limit} \times \text{Gas Price}
\end{align*}

\textbf{Refund}: If transaction uses less gas than limit, difference is refunded:
\[
\text{Refund} = (\text{Gas Limit} - \text{Gas Used}) \times \text{Gas Price}
\]
\end{keybox}

\subsection{Gas Costs by Operation}

\begin{center}
\begin{tabular}{|l|r|l|}
\hline
\textbf{Operation} & \textbf{Gas Cost} & \textbf{Category} \\
\hline
Addition (ADD) & 3 & Arithmetic \\
Multiplication (MUL) & 5 & Arithmetic \\
Division (DIV) & 5 & Arithmetic \\
Storage Write (SSTORE) & 20,000 / 5,000 & Storage \\
Storage Read (SLOAD) & 200 & Storage \\
Keccak-256 Hash & 30 + 6/word & Crypto \\
Call to Contract & 700 & Call \\
Create Contract & 32,000 & Create \\
Transaction Base & 21,000 & Base \\
\hline
\end{tabular}
\end{center}

\subsection{EIP-1559: Gas Fee Reform}

Introduced in August 2021 (London hard fork):

\textbf{New Fee Structure}:
\[
\text{Total Fee} = \text{Gas Used} \times (\text{Base Fee} + \text{Priority Fee})
\]

\begin{itemize}
  \item \textbf{Base Fee}: Algorithmically determined, burned (not given to miners)
  \item \textbf{Priority Fee} (Tip): User sets, goes to validators
  \item \textbf{Max Fee}: User's maximum willingness to pay
\end{itemize}

\textbf{Base Fee Adjustment}:
\[
\text{New Base Fee} = \text{Old Base Fee} \times \left(1 + \frac{8(\text{Gas Used} - \text{Target})}{(\text{Target} \times 8)}\right)
\]

Increases by up to 12.5\% if block is full, decreases if empty.

\begin{examplebox}{Gas Calculation Example}
User submits transaction with:
\begin{itemize}
  \item Gas Limit: 100,000
  \item Max Fee: 200 gwei
  \item Priority Fee: 2 gwei
\end{itemize}

If base fee is 150 gwei and transaction uses 80,000 gas:

\begin{align*}
\text{Actual Fee per Gas} &= 150 + 2 = 152 \text{ gwei} \\
\text{Total Cost} &= 80,000 \times 152 = 12,160,000 \text{ gwei} = 0.01216 \text{ ETH} \\
\text{Burned} &= 80,000 \times 150 = 12,000,000 \text{ gwei} \\
\text{To Validator} &= 80,000 \times 2 = 160,000 \text{ gwei}
\end{align*}
\end{examplebox}

\section{Solidity Programming Language}

\subsection{Overview}

Solidity is a high-level, statically-typed language for writing smart contracts. It compiles to EVM bytecode.

\textbf{Characteristics}:
\begin{itemize}
  \item Syntax similar to JavaScript/C++
  \item Strongly typed
  \item Supports inheritance, libraries, user-defined types
  \item Version specified with pragma directive
\end{itemize}

\subsection{Basic Contract Structure}

\begin{examplebox}{Minimal Contract}
\begin{lstlisting}[language=Solidity]
// SPDX-License-Identifier: MIT
pragma solidity ^0.8.0;

contract SimpleStorage {
    // State variable stored on blockchain
    uint256 public storedData;

    // Constructor called once at deployment
    constructor(uint256 initialValue) {
        storedData = initialValue;
    }

    // Function to update storage
    function set(uint256 newValue) public {
        storedData = newValue;
    }

    // View function (read-only, no gas when called externally)
    function get() public view returns (uint256) {
        return storedData;
    }
}
\end{lstlisting}
\end{examplebox}

\subsection{Data Types}

\subsubsection{Value Types}

\begin{itemize}
  \item \textbf{Boolean}: \texttt{bool} (true/false)
  \item \textbf{Integers}:
    \begin{itemize}
      \item Unsigned: \texttt{uint8} to \texttt{uint256} (steps of 8)
      \item Signed: \texttt{int8} to \texttt{int256}
      \item Default \texttt{uint} = \texttt{uint256}
    \end{itemize}
  \item \textbf{Address}: \texttt{address} (20 bytes), \texttt{address payable}
  \item \textbf{Bytes}: \texttt{bytes1} to \texttt{bytes32}, \texttt{bytes} (dynamic)
  \item \textbf{String}: \texttt{string} (dynamic UTF-8)
  \item \textbf{Enums}: Custom enumerated types
\end{itemize}

\subsubsection{Reference Types}

\begin{itemize}
  \item \textbf{Arrays}: \texttt{uint[]} (dynamic), \texttt{uint[10]} (fixed)
  \item \textbf{Structs}: Custom composite types
  \item \textbf{Mappings}: \texttt{mapping(address => uint)}
\end{itemize}

\textbf{Data Location} (must specify for reference types):
\begin{itemize}
  \item \texttt{storage}: Persistent, expensive
  \item \texttt{memory}: Temporary, cheaper
  \item \texttt{calldata}: Read-only, function parameters
\end{itemize}

\begin{examplebox}{Advanced Types Example}
\begin{lstlisting}[language=Solidity]
pragma solidity ^0.8.0;

contract AdvancedTypes {
    // Struct definition
    struct Person {
        string name;
        uint age;
        address wallet;
    }

    // Mapping from address to Person
    mapping(address => Person) public people;

    // Dynamic array
    address[] public allAddresses;

    // Enum
    enum Status { Pending, Active, Closed }
    Status public currentStatus;

    function addPerson(string memory _name, uint _age) public {
        people[msg.sender] = Person(_name, _age, msg.sender);
        allAddresses.push(msg.sender);
        currentStatus = Status.Active;
    }

    function getPerson(address _addr) public view
        returns (string memory, uint, address) {
        Person memory p = people[_addr];
        return (p.name, p.age, p.wallet);
    }
}
\end{lstlisting}
\end{examplebox}

\subsection{Function Modifiers}

\subsubsection{Visibility}

\begin{itemize}
  \item \texttt{public}: Callable internally and externally
  \item \texttt{external}: Only callable externally (more gas efficient)
  \item \texttt{internal}: Only within contract and derived contracts
  \item \texttt{private}: Only within current contract
\end{itemize}

\subsubsection{State Mutability}

\begin{itemize}
  \item \texttt{view}: Reads state, doesn't modify (no gas when called externally)
  \item \texttt{pure}: Doesn't read or modify state (no gas when called externally)
  \item \texttt{payable}: Can receive Ether
  \item (none): Can modify state (costs gas)
\end{itemize}

\subsection{Special Variables and Functions}

\begin{center}
\begin{tabular}{|l|p{7cm}|}
\hline
\textbf{Variable} & \textbf{Description} \\
\hline
\texttt{msg.sender} & Address of immediate caller \\
\texttt{msg.value} & Amount of wei sent with transaction \\
\texttt{msg.data} & Complete calldata \\
\texttt{block.timestamp} & Current block timestamp (Unix epoch) \\
\texttt{block.number} & Current block number \\
\texttt{tx.origin} & Original transaction sender (EOA) \\
\texttt{address(this)} & Current contract address \\
\texttt{address.balance} & Balance of address in wei \\
\hline
\end{tabular}
\end{center}

\subsection{Events}

Events provide logging facility for DApps to listen to state changes.

\begin{examplebox}{Event Usage}
\begin{lstlisting}[language=Solidity]
pragma solidity ^0.8.0;

contract EventExample {
    // Event declaration (indexed parameters can be filtered)
    event Transfer(
        address indexed from,
        address indexed to,
        uint256 value
    );

    mapping(address => uint256) public balances;

    function transfer(address to, uint256 amount) public {
        require(balances[msg.sender] >= amount, "Insufficient balance");

        balances[msg.sender] -= amount;
        balances[to] += amount;

        // Emit event
        emit Transfer(msg.sender, to, amount);
    }
}
\end{lstlisting}

\textbf{Event Storage}: Events are stored in transaction logs (not in contract storage), making them much cheaper for recording information.
\end{examplebox}

\section{Smart Contract Lifecycle}

\subsection{Development Workflow}

\begin{enumerate}
  \item \textbf{Write}: Develop Solidity code
  \item \textbf{Compile}: Convert to bytecode using \texttt{solc} compiler
  \item \textbf{Test}: Unit tests on local blockchain (Hardhat, Ganache)
  \item \textbf{Deploy}: Send creation transaction to network
  \item \textbf{Verify}: Publish source code on Etherscan
  \item \textbf{Interact}: Call functions via transactions
  \item \textbf{Upgrade} (if designed): Use proxy patterns for upgradeability
\end{enumerate}

\subsection{Contract Deployment}

Deployment creates a special transaction with:
\begin{itemize}
  \item \texttt{to} address: empty (null)
  \item \texttt{data}: Contract bytecode + constructor parameters
  \item \texttt{gas}: Sufficient for deployment
\end{itemize}

\textbf{Contract Address Calculation}:
\[
\text{Address} = \text{Keccak256}(\text{RLP}(\text{sender\_address}, \text{nonce}))[:20]
\]

\begin{examplebox}{Deployment with Hardhat}
\begin{lstlisting}[language=JavaScript]
// scripts/deploy.js
const hre = require("hardhat");

async function main() {
  // Get contract factory
  const SimpleStorage = await hre.ethers.getContractFactory(
    "SimpleStorage"
  );

  // Deploy with constructor argument
  const contract = await SimpleStorage.deploy(42);

  await contract.deployed();

  console.log("Contract deployed to:", contract.address);
  console.log("Initial value:", await contract.get());
}

main().catch((error) => {
  console.error(error);
  process.exitCode = 1;
});
\end{lstlisting}
\end{examplebox}

\subsection{Contract Interaction}

\subsubsection{Reading State (Calls)}

\begin{itemize}
  \item Execute locally on node
  \item Don't create transactions
  \item Free (no gas cost)
  \item Return values immediately
  \item Used for \texttt{view} and \texttt{pure} functions
\end{itemize}

\subsubsection{Modifying State (Transactions)}

\begin{itemize}
  \item Broadcast to network
  \item Require gas payment
  \item Mined into block (confirmation delay)
  \item Generate transaction receipt
  \item Return transaction hash, not function return value
\end{itemize}

\section{Common Design Patterns}

\subsection{Access Control}

\begin{examplebox}{Ownable Pattern}
\begin{lstlisting}[language=Solidity]
pragma solidity ^0.8.0;

contract Ownable {
    address public owner;

    event OwnershipTransferred(
        address indexed previousOwner,
        address indexed newOwner
    );

    constructor() {
        owner = msg.sender;
    }

    modifier onlyOwner() {
        require(msg.sender == owner, "Not the owner");
        _;
    }

    function transferOwnership(address newOwner) public onlyOwner {
        require(newOwner != address(0), "Invalid address");
        emit OwnershipTransferred(owner, newOwner);
        owner = newOwner;
    }
}

contract MyContract is Ownable {
    uint256 public value;

    // Only owner can call this
    function setValue(uint256 _value) public onlyOwner {
        value = _value;
    }
}
\end{lstlisting}
\end{examplebox}

\subsection{Checks-Effects-Interactions}

Best practice to prevent reentrancy attacks:

\begin{enumerate}
  \item \textbf{Checks}: Validate conditions
  \item \textbf{Effects}: Update state
  \item \textbf{Interactions}: Call external contracts
\end{enumerate}

\begin{warningbox}{Reentrancy Vulnerability}
\textbf{Vulnerable Code}:
\begin{lstlisting}[language=Solidity]
function withdraw(uint amount) public {
    require(balances[msg.sender] >= amount);
    // BAD: External call before state update
    msg.sender.call{value: amount}("");
    balances[msg.sender] -= amount;
}
\end{lstlisting}

\textbf{Secure Code}:
\begin{lstlisting}[language=Solidity]
function withdraw(uint amount) public {
    require(balances[msg.sender] >= amount);
    // GOOD: Update state first
    balances[msg.sender] -= amount;
    msg.sender.call{value: amount}("");
}
\end{lstlisting}

Famous example: The DAO hack (2016) exploited reentrancy, leading to \$60M loss and Ethereum/Ethereum Classic fork.
\end{warningbox}

\section{Study Questions}

\subsection{Conceptual Questions}

\begin{enumerate}
  \item Explain the difference between \texttt{memory} and \texttt{storage} in Solidity.
  \item Why is storage more expensive than memory or calldata?
  \item What happens if a transaction runs out of gas during execution?
  \item How does EIP-1559 improve on the legacy gas pricing mechanism?
  \item Describe the purpose of events and why they're cheaper than storage.
\end{enumerate}

\subsection{Design Problems}

\begin{enumerate}
  \item Design a voting contract where each address can vote once on proposals.
  \item Implement a simple escrow contract that holds funds until both parties agree.
  \item Create a contract that tracks employee work hours and calculates payment.
  \item Build a crowdfunding contract with refund mechanism if goal isn't met.
\end{enumerate}

\subsection{Security Analysis}

\begin{enumerate}
  \item Identify the vulnerability in a given contract code snippet.
  \item Explain how the checks-effects-interactions pattern prevents reentrancy.
  \item Why should you avoid using \texttt{tx.origin} for authentication?
  \item What are the risks of using \texttt{block.timestamp} for critical logic?
\end{enumerate}

\section{Further Reading}

\begin{itemize}
  \item Ethereum Yellowpaper (Formal specification)
  \item Solidity Documentation: \url{https://docs.soliditylang.org}
  \item OpenZeppelin Contracts: Secure, audited implementations
  \item Consensys Smart Contract Best Practices
  \item Trail of Bits: Building Secure Contracts
  \item Ethereum Improvement Proposals (EIPs)
\end{itemize}

\end{document}
